\chapter{METODE PENELITIAN}

\section{Kerangka Berpikir}
Berikut ini gambar Kerangka Berpikir yang digunakan dalam penelitian:
\begin{figure}[h]
    \centering
    \includegraphics[width=10cm, height=13cm]{/home/veinmahzy/skrispi/images/bab_3/Kerangka Bepikir.png}
    \caption{Kerangka Berpikir}
    \label{fig:kerangka-berfikir}
\end{figure}

Penelitian ini mengusulkan pengembangan sistem klasifikasi otomatis berbasis kecerdasan buatan untuk mengategorikan laporan masyarakat ke dalam delapan kelas spesifik, meliputi infrastruktur (jalanan rusak, kerusakan trotoar, halte, dan pohon tumbang), gangguan lingkungan (vandalisme dan sampah), serta insiden lalu lintas. 
Siklus pengembangan dimulai dengan optimasi performa model melalui fase pelatihan yang rigid, diikuti oleh perancangan arsitektur perangkat lunak yang mengintegrasikan \textit{AI Service} ke dalam ekosistem web dengan struktur basis data yang terukur. 
Setelah tahap implementasi (\textit{development}), sistem divalidasi melalui survei pengguna untuk mengevaluasi aspek fungsionalitas dan pengalaman pengguna. Tahap akhir penelitian melibatkan proses revisi iteratif guna memastikan reliabilitas dan efektivitas sistem sebelum diimplementasikan secara luas oleh masyarakat dan instansi pemerintah terkait.

\section{Metode Penelitian}
Penelitian ini berfokus pada integrasi arsitektur \textit{Deep Learning} ke dalam platform pelaporan publik untuk menghasilkan sistem klasifikasi otomatis dengan tingkat presisi yang tinggi. Metodologi penelitian dikonstruksi secara sistematis dalam dua fase utama: optimalisasi model berbasis kecerdasan buatan (\textit{Artificial Intelligence}) dan perancangan infrastruktur sistem berbasis web. Dalam pengembangan model, dilakukan studi komparatif untuk mengevaluasi efektivitas arsitektur \textit{Residual Network} (ResNet101) sebagai representasi model \textit{Convolutional Neural Network} (CNN) konvensional terhadap arsitektur \textit{Vision Transformer} (ViT) yang berbasis mekanisme \textit{attention}. Integrasi alur kerja antara komponen model dan sistem manajemen data tersebut dipaparkan secara komprehensif pada Gambar \ref{fig:kerangka-berfikir}.

\subsection{Pengembangan Sistem \textit{Artificial Intelligence} (AI)}
Alur pengembangan kecerdasan buatan dalam penelitian ini dipaparkan secara sistematis pada Gambar \ref{fig:Alur Eskperimen AI}. Tahap inisiasi dimulai dengan pengumpulan, akuisisi, serta prapemrosesan (\textit{preprocessing}) dataset guna memastikan integritas data. Selanjutnya, penelitian ini menerapkan teknik ekstraksi fitur multimodal untuk mentransformasi data teks dan gambar menjadi vektor \textit{embedding}. 

Peneliti melakukan eksperimen komprehensif melalui kombinasi arsitektur konvensional dan \textit{transformer}, yang mencakup ResNet101, DINOv2, FastText, IndoBERT, dan E5. Representasi fitur yang dihasilkan dari kombinasi model-model tersebut kemudian diintegrasikan melalui mekanisme \textit{feature concatenation}. Terakhir, evaluasi performa dilakukan secara rigid menggunakan dataset uji dengan metrik akurasi sebagai parameter utama untuk menentukan konfigurasi model terbaik.
\begin{figure}[h]
    \centering
    \includegraphics[width=\textwidth]{/home/veinmahzy/skrispi/images/bab_3/Artifical Intelegence Alur Eskperimen.png}
    \caption{Alur Eskperimen AI}
    \label{fig:Alur Eskperimen AI}
\end{figure}

\subsubsection{Dataset}
Prosedur akuisisi data dalam penelitian ini dilakukan melalui pendekatan hibrida yang mengombinasikan sumber data terbuka (\textit{open-source}) dan observasi mandiri. Pertama, dataset sekunder diperoleh dari platform Kaggle dengan fokus spesifik pada domain deteksi kerusakan infrastruktur jalan. Kedua, dilakukan ekstraksi data primer secara mandiri menggunakan teknik \textit{web scraping} pada portal resmi pengaduan masyarakat (\texttt{crm.jakarta.go.id}). 

Proses \textit{scraping} ini berhasil menghimpun sebanyak 15.885 entri laporan publik yang tercatat sepanjang tahun 2025 yang bersifat publik. Struktur dataset mentah mencakup metadata heterogen yang terdiri dari:
\begin{itemize}
    \item Identitas unik laporan (\textit{Report ID}).
    \item Gambar(\textit{Supporting images}) sebagai input fitur visual.
    \item Deskripsi tekstual laporan sebagai input fitur \textit{Natural Language Processing} (NLP).
    \item Metrik waktu disposisi kelurahan ke instansi terkait.
    \item Label instansi penindaklanjuti.
\end{itemize}

laporan tersebut.
Dataset tersebut kemudian diolah dan dibersihkan untuk
mempersiapkan training model dengan target klasifikasi untuk mempersiapkan \textit{training} model dengan target klasifikasi untuk delapan permasalahan laporan yaitu jalanan rusak, kecelakaan, kerusakan trotoar, vandalisme, sampah, kerusakan halte, kecelakaan lalu lintas, dan pohon tumbang

\subsubsection{Metrik Evaluasi}

Pemilihan akurasi sebagai metrik evaluasi utama didasarkan pada asumsi \textit{uniform importance}, di mana setiap kategori laporan masyarakat dianggap memiliki bobot kepentingan yang setara tanpa memandang instansi terkait. Akurasi dipilih karena kemampuannya dalam memberikan representasi performa global model secara intuitif. Metrik ini mengukur rasio antara jumlah prediksi yang benar terhadap keseluruhan total observasi dalam dataset.

Secara matematis, akurasi didefinisikan berdasarkan komponen \textit{confusion matrix} sebagai berikut:

\begin{equation}
    \label{eq:accuracy_full}
    \text{Accuracy} = \frac{TP + TN}{TP + TN + FP + FN}
\end{equation}

Di mana $TP$ (\textit{True Positive}) dan $TN$ (\textit{True Negative}) merepresentasikan jumlah klasifikasi yang tepat, sedangkan $FP$ (\textit{False Positive}) dan $FN$ (\textit{False Negative}) merepresentasikan kesalahan klasifikasi. Dalam konteks penelitian ini, formulasi tersebut dapat disederhanakan menjadi:

\begin{equation}
    \label{eq:accuracy_simple}
    acc = \frac{P}{N}
\end{equation}

di mana $P$ adalah jumlah laporan yang diprediksi dengan benar secara keseluruhan dan $N$ adalah total populasi data. Penggunaan metrik ini memastikan bahwa evaluasi berfokus pada efektivitas sistem secara holistik dalam mengklasifikasikan setiap laporan yang masuk secara adil.

\subsubsection{Pengaturan Eksperimen}
Pada Tabel \ref{tab:model-all} adalah \textit{weight checkpoint} untuk model teks dan gambar, semuanya merupakan \textit{open source model} yang dapat diakses melalui Hugging Face.

\begin{table}[htbp]
    \centering
    \caption{Detail Model Deep Learning}
    \label{tab:model-all}
    \begin{tabular}{lll}
        \toprule
        \textbf{Jenis} & \textbf{Nama Model} & \textbf{Model Checkpoint} \\
        \midrule
        \multirow{3}{*}{Gambar} 
        & ResNet101 & \texttt{microsoft/resnet-101} \\
        & ViT-B/16 & \texttt{google/vit-base-patch16-224-in21k} \\
        & DINOv2 & \texttt{facebook/dinov2-large} \\
        \midrule
        \multirow{3}{*}{Teks}
        & FastText & \texttt{facebook/fasttext-id-vectors} \\
        & IndoBERT & \texttt{cahya/bert-base-indonesian-522M} \\
        & Multilingual E5 & \texttt{intfloat/multilingual-e5-large-instruct} \\
        \bottomrule
    \end{tabular}
\end{table}